\documentclass[11pt]{article}
\usepackage[margin=1in]{geometry}
\usepackage{amsmath,amssymb,amsthm,mathtools}
\usepackage{graphicx}
\usepackage[hidelinks]{hyperref}

\newtheorem{theorem}{Theorem}
\newtheorem{lemma}[theorem]{Lemma}
\newcommand{\C}{\mathbb C}
\newcommand{\Rp}{\mathfrak R}
\newcommand{\Sp}{\mathcal S}
\newcommand{\had}{\mathbin{*}}
\setlength{\emergencystretch}{2em}

\title{The predual of bounded Schatten Schur multipliers\\
is a Banach algebra under matrix multiplication}
\author{Autonomous mathematical discovery run \texttt{fa\_banach\_001}}
\date{17 August 2026}

\begin{document}
\maketitle

\noindent\textbf{Status.}
Candidate full affirmative answer, likely valid pending expert review.

\section{Source question}

Let $I$ be a set and $1\leq p<\infty$.  Arhancet defines
$\Rp_p^I$ as the natural predual of the bounded Schur multipliers on
$\Sp_p^I$.  Concretely, $\Rp_p^I$ is the image of
\[
 \Sp_p^I\widehat\otimes_\pi\Sp_{p'}^I\longrightarrow \Sp_1^I,
 \qquad A\otimes B\longmapsto A\had B,
\]
with its quotient norm.  Remark~3.8 asks whether $\Rp_p^I$, equipped with
ordinary matrix multiplication, is a Banach algebra
\cite[Remark~3.8, PDF p.~18]{Arhancet}.

\begin{center}
\includegraphics[width=0.96\textwidth]{figures/open_question.png}
\end{center}

We answer this affirmatively, with multiplication constant one.  The point is
that one only needs to amplify the very special tensors arising from two
Hadamard factorizations.  A rank-one decomposition controls those tensors
directly, even though the general Banach-space amplification used in the
source is false.

\section{Preliminaries from the source}

Write $A\had B=(a_{ij}b_{ij})$.  The quotient description gives
\begin{equation}\label{eq:norm}
 \|C\|_{\Rp_p^I}
 =\inf\left\{
   \sum_n\|A_n\|_p\|B_n\|_{p'}:
   C=\sum_n A_n\had B_n
 \right\}.
\end{equation}
The series in this formula is absolutely convergent in the projective
tensor-product sense.  The canonical inclusion $\Rp_p^I\hookrightarrow
\Sp_1^I$ is contractive.  For $1<p<\infty$, the identity on matrices is an
isometric identification
\begin{equation}\label{eq:duality}
 \Rp_p^I=\Rp_{p'}^I;
\end{equation}
these facts are recorded in Section~2 of \cite{Arhancet}.

\section{Idea of the proof}

Start with two elementary quotient representatives $A\had B$ and
$X\had Y$.  Their ordinary matrix product is a sum over the middle index
$k$.  For each $k$, put the corresponding column--row outer products into
rank-one matrices $P_k$ and $Q_k$; the $k$th summand is then exactly
$P_k\had Q_k$.  When $p\geq2$, the inclusion
$\Sp_{p'}^I\subseteq\Sp_2^I$ turns the total projective cost of these
rank-one factorizations into a Cauchy--Schwarz estimate with constant one.
Summing this elementary estimate over arbitrary quotient representations
proves submultiplicativity.  The source isometry
$\Rp_p^I=\Rp_{p'}^I$ transfers the result to $1<p<2$, while the source's
completely bounded endpoint theorem handles $p=1$.

\section{The elementary product estimate}

\begin{lemma}\label{lem:elementary}
Let $2\leq p<\infty$.  For $A,X\in\Sp_p^I$ and
$B,Y\in\Sp_{p'}^I$,
\[
 (A\had B)(X\had Y)\in\Rp_p^I
\]
and
\begin{equation}\label{eq:elementary}
 \|(A\had B)(X\had Y)\|_{\Rp_p^I}
 \leq \|A\|_p\|B\|_{p'}\|X\|_p\|Y\|_{p'}.
\end{equation}
\end{lemma}

\begin{proof}
For $k\in I$, define the rank-one matrices
\[
 P_k=(a_{ik}x_{kj})_{i,j\in I},
 \qquad Q_k=(b_{ik}y_{kj})_{i,j\in I}.
\]
Thus $P_k$ is the outer product of column $k$ of $A$ and row $k$ of
$X$, and similarly for $Q_k$.  Every Schatten norm of a rank-one matrix is
the product of the two Euclidean norms, so, with
\[
 \alpha_k=\|Ae_k\|_2,\quad \gamma_k=\|X^*e_k\|_2,
 \quad \beta_k=\|Be_k\|_2,\quad \delta_k=\|Y^*e_k\|_2,
\]
we have
\[
 \|P_k\|_p=\alpha_k\gamma_k,
 \qquad \|Q_k\|_{p'}=\beta_k\delta_k.
\]
Because $p'\leq2$, the inclusion $\Sp_{p'}^I\subseteq\Sp_2^I$ is
contractive.  Hence Cauchy--Schwarz gives
\begin{align*}
 \sum_k\|P_k\|_p\|Q_k\|_{p'}
 &=\sum_k\alpha_k\gamma_k\beta_k\delta_k\\
 &\leq \big(\sup_k\alpha_k\gamma_k\big)
          \sum_k\beta_k\delta_k\\
 &\leq \|A\|_\infty\|X\|_\infty
          \|B\|_2\|Y\|_2\\
 &\leq \|A\|_p\|X\|_p\|B\|_{p'}\|Y\|_{p'}.
\end{align*}
Therefore $\sum_kP_k\otimes Q_k$ is absolutely convergent in
$\Sp_p^I\widehat\otimes_\pi\Sp_{p'}^I$.  Applying the Hadamard-product
quotient map and comparing entries yields
\[
 \sum_kP_k\had Q_k
 =\left(\sum_k a_{ik}b_{ik}x_{kj}y_{kj}\right)_{i,j}
 =(A\had B)(X\had Y).
\]
The norm formula \eqref{eq:norm} now gives \eqref{eq:elementary}.
\end{proof}

\section{Full answer}

\begin{theorem}\label{thm:main}
For every set $I$ and every $1\leq p<\infty$, $\Rp_p^I$ is a
contractive Banach algebra under ordinary matrix multiplication.  Explicitly,
for all $C,D\in\Rp_p^I$,
\begin{equation}\label{eq:main}
 CD\in\Rp_p^I,
 \qquad \|CD\|_{\Rp_p^I}\leq
 \|C\|_{\Rp_p^I}\|D\|_{\Rp_p^I}.
\end{equation}
Moreover, for every $i,j\in I$, the net
$\sum_{k\in J}c_{ik}d_{kj}$ over finite subsets $J\subset I$ converges
absolutely to $(CD)_{ij}$.
\end{theorem}

\begin{proof}
First suppose $2\leq p<\infty$.  Choose representations
\[
 C=\sum_nA_n\had B_n,
 \qquad D=\sum_mX_m\had Y_m
\]
as in \eqref{eq:norm}.  Lemma~\ref{lem:elementary} shows that the double
series
\[
 \sum_{n,m}(A_n\had B_n)(X_m\had Y_m)
\]
converges absolutely in $\Rp_p^I$, with norm at most
\[
 \left(\sum_n\|A_n\|_p\|B_n\|_{p'}\right)
 \left(\sum_m\|X_m\|_p\|Y_m\|_{p'}\right).
\]
The inclusion $\Rp_p^I\hookrightarrow\Sp_1^I$ is contractive, and
$\Sp_1^I$ is a Banach algebra under operator composition.  Consequently the
sum is exactly $CD$.  Taking infima over the two representations proves
\eqref{eq:main}.

If $1<p<2$, then $p'>2$ and the isometry \eqref{eq:duality} transfers the
already proved estimate from $\Rp_{p'}^I$ to $\Rp_p^I$.  For $p=1$, the
source records $\Rp_1^I=\Rp_{1,cb}^I$ isometrically, and its Theorem~3.7
already proves the contractive multiplication estimate.  This completes all
values of $p$.

Finally, since $C,D\in\Sp_1^I\subset\Sp_2^I$, every row of $C$ and every
column of $D$ lies in $\ell_2(I)$.  Thus
\[
 \sum_k|c_{ik}d_{kj}|
 \leq \left(\sum_k|c_{ik}|^2\right)^{1/2}
      \left(\sum_k|d_{kj}|^2\right)^{1/2}<\infty,
\]
which proves the stated coordinate convergence.
\end{proof}

\section{Scope, verification, and novelty}

Theorem~\ref{thm:main} fully answers the first question in Remark~3.8 for
all exponents and all index sets.  It does not settle the companion question
whether every bounded Schur multiplier on $\Sp_p$ is completely bounded.
The proof deliberately avoids that stronger conjecture.

The accompanying script in \texttt{code/verify\_factorization.py} verifies,
on random finite matrices and several
exponents, the rank-one identity and every numerical inequality in
Lemma~\ref{lem:elementary}.  These computations are regression checks; the
proof itself is exact and dimension-free.

A bounded novelty search, refreshed on 21 August 2026, covered the four run indexes,
exact title and arXiv-id searches, exact phrases from Remark~3.8, and later
work on the complete bounds of Schatten Schur multipliers.  No later answer
to the $\Rp_p$ matrix-product question was found.  The later literature found
continues to study the distinct companion equality
$\mathfrak M_p=\mathfrak M_{p,cb}$.

\paragraph{Human review recommendation.}
Check the four-factor rank-one decomposition in Lemma~\ref{lem:elementary},
especially the use of $\Sp_{p'}\subset\Sp_2$ for $p\geq2$, and then check
that the projective representations pass to the operator product through the
contractive inclusion $\Rp_p^I\hookrightarrow\Sp_1^I$.

\begin{thebibliography}{9}
\raggedright
\bibitem{Arhancet}
C. Arhancet,
\emph{Noncommutative Fig\`a-Talamanca--Herz algebras for Schur multipliers},
Integral Equations Operator Theory \textbf{70} (2011), 485--510;
arXiv:0911.4304v2.
\end{thebibliography}

\end{document}
